% force the table of contents to include our abstract
\cleardoublepage
\addcontentsline{toc}{section}{Abstract}\chapter*{Abstract}

Sparse matrix-matrix multiplication (SpGEMM) is a core kernel in graph analytics, scientific computing and machine learning, and how fast it runs depends heavily on how the input matrices are stored. This thesis asks how the best pairing of storage formats changes with the size and structure of the input, and what the hardware is doing when it does. This thesis argues that, for sparse matrix-matrix multiplication using the kernels I chose and on the Kronecker graph data I tested on, choosing among storage formats is a simple decision. The dataflow of the format pair matters by orders of magnitude, and within the good pairs a fixed choice is within a few percent of the best. However, where the best pair does change follows quantities such as Number of Non-zeroes(nnz) or size, known before the multiply runs.

This thesis contributes the following: A benchmarking library that records timing and hardware counters (cache misses, branch mispredictions, page faults) around marked regions of code. Floating point and memory operation formulae for all sixteen pairings of four formats, which matched measured counts for fourteen of them, but operation counts turned out to be a poor predictor of runtime. A negative result in implementing a pipeline that picks a format for every block of a matrix based on the previous formulae and runs the multiply through a BLIS-style loop structure did not beat a single global format, because dispatch overhead dominated. A benchmarking sweep of every format pair over 19 different Kronecker matrices, which used my benchmarking library to collect hardware counter data in addition to timing. This sweep demonstrated that inner-product dataflows are catastrophically slow, and that at low K a Gustavson dataflow using \emph{CSRxCSR} is the best pairing, before giving way to outer product dataflows like \emph{CSCxCSR} and \emph{CSCxCOO}, with \emph{CSCxCOO} winning out of the 2 when nnz/row is very low. And finally another partly negative, results of tiling on the same 19 matrices which found little benefit or in fact major performance loss. However, some matrices show minor promising trends where the relative cost of blocking fell as the matrix got larger. The thesis also includes an earlier FPGA design for streaming multi-format sparse matrix-vector multiplication, which motivated these questions.

These conclusions are limited to the specific kernel implementations I used and the Kronecker data sets I tested on.
